% GENERATED FILE. Edit canonical lessons, then run npm run generate:books.
% canonical-source-hash: fnv1a64:28f533068f81b06b
% canonical-lessons: HI-C424-dhaba, HI-C424-halvai, HI-C424-thela, HI-C424-pherivala, HI-C424-pan

\chapter{Street Food and Sellers}
\label{ch:hi-street-food-and-sellers}
\hlchaptermodality{424}

\begin{chapteropening}
\textbf{By the end of this chapter:} \emph{I can eat at a roadside dhaba, buy sweets and paan, and talk about the hawkers in the lane.}

Complete the last lesson of chapter 424: I can eat at a roadside dhaba, buy sweets and paan, and talk about the hawkers in the lane.
\end{chapteropening}

\section[\texorpdfstring{\hi{ढाबा} (ḍhābā)}{ढाबा (ḍhābā)}]{\hi{ढाबा} (ḍhābā) — a roadside eatery, a dhaba}

\label{lesson:HI-C424-dhaba}



\begin{quote}
Before the new one: say the Hindi for to growl, then the Hindi for to bark.
\end{quote}

\subsection*{You'll want to know: \hi{ढाबा}}
\textbf{\hi{ढाबा}} — \emph{ḍhābā} — \textquotedblleft{}a roadside eatery, a dhaba\textquotedblright{}.

\begin{quote}
\hi{राजमार्ग} \hi{पर} \hi{हम} \hi{एक} \hi{ढाबे} \hi{पर} \hi{रुके।}
\end{quote}

\begin{quote}
\emph{rājmārg par ham ek ḍhābe par ruke.}
\end{quote}

\begin{quote}
On the highway we stopped at a dhaba.
\end{quote}

Masculine.

\begin{cousinweb}
\textbf{\hi{ढाबा}} (\emph{ḍhābā}) goes back to \emph{*ḍhapp}, \textquotedblleft{}cover\textquotedblright{}, a form CDIAL reconstructs (CDIAL 5579), which lists Hindi \emph{ḍhābā} as \textquotedblleft{}thatched roofing\textquotedblright{}. The roadside eatery is likely named for its thatched shelter, though CDIAL does not give that sense.
\end{cousinweb}

\begin{grammarlens}[title={What you've built}]
One more word for getting around.
\end{grammarlens}

\subsection*{Your turn}
\begin{itemize}
\raggedright
  \item \emph{Say it:} \emph{ḍhābā}
  \item \emph{Say it:} \emph{ḍhābā}, once more
  \item \emph{Say it:} \emph{rājmārg par ham ek ḍhābe par ruke.} — On the highway we stopped at a dhaba.
  \item \emph{Recall:} say the Hindi for to growl, then the Hindi for to bark, then say \emph{ḍhābā} again
\end{itemize}

\begin{tcolorbox}[breakable,colback=teal!4,colframe=teal!35!black,title={Before you move on}]
What does \hi{ढाबा} mean? (\textquotedblleft{}A roadside eatery, a dhaba\textquotedblright{}.) Say it once more.
\end{tcolorbox}
\section[\texorpdfstring{\hi{हलवाई} (halvāī)}{हलवाई (halvāī)}]{\hi{हलवाई} (halvāī) — a sweet-maker; a sweet shop}

\label{lesson:HI-C424-halvai}



\begin{quote}
Before the new one: say the Hindi for to bark, then the Hindi for a roadside eatery.
\end{quote}

\subsection*{You'll want to know: \hi{हलवाई}}
\textbf{\hi{हलवाई}} — \emph{halvāī} — \textquotedblleft{}a sweet-maker; a sweet shop\textquotedblright{}.

\begin{quote}
\hi{हलवाई} \hi{की} \hi{दुकान} \hi{से} \hi{मैंने} \hi{लड्डू} \hi{ख़रीदे।}
\end{quote}

\begin{quote}
\emph{halvāī kī dukān se maiṁne laḍḍū kharīde.}
\end{quote}

\begin{quote}
I bought laddus from the sweet shop.
\end{quote}

Masculine.

\begin{cousinweb}
\textbf{\hi{हलवाई}} (\emph{halvāī}) is built on \hi{हलवा}, which is usually traced to Arabic \emph{ḥalwā}, \textquotedblleft{}a sweetmeat\textquotedblright{}, through Persian; the ending -āī makes the person who makes it.
\end{cousinweb}

\begin{grammarlens}[title={What you've built}]
One more word for getting around.
\end{grammarlens}

\subsection*{Your turn}
\begin{itemize}
\raggedright
  \item \emph{Say it:} \emph{halvāī}
  \item \emph{Say it:} \emph{halvāī}, once more
  \item \emph{Say it:} \emph{halvāī kī dukān se maiṁne laḍḍū kharīde.} — I bought laddus from the sweet shop.
  \item \emph{Recall:} say the Hindi for to bark, then the Hindi for a roadside eatery, then say \emph{halvāī} again
\end{itemize}

\begin{tcolorbox}[breakable,colback=teal!4,colframe=teal!35!black,title={Before you move on}]
What does \hi{हलवाई} mean? (\textquotedblleft{}A sweet-maker; a sweet shop\textquotedblright{}.) Say it once more.
\end{tcolorbox}
\section[\texorpdfstring{\hi{ठेला} (ṭhelā)}{ठेला (ṭhelā)}]{\hi{ठेला} (ṭhelā) — a handcart}

\label{lesson:HI-C424-thela}



\begin{quote}
Before the new one: say the Hindi for a roadside eatery, then the Hindi for a sweet-maker.
\end{quote}

\subsection*{You'll want to know: \hi{ठेला}}
\textbf{\hi{ठेला}} — \emph{ṭhelā} — \textquotedblleft{}a handcart\textquotedblright{}.

\begin{quote}
\hi{एक} \hi{आदमी} \hi{ठेले} \hi{पर} \hi{सब्ज़ी} \hi{बेच} \hi{रहा} \hi{था।}
\end{quote}

\begin{quote}
\emph{ek ādmī ṭhele par sabzī bec rahā thā.}
\end{quote}

\begin{quote}
A man was selling vegetables from a handcart.
\end{quote}

Masculine.

\begin{cousinweb}
\textbf{\hi{ठेला}} (\emph{ṭhelā}) goes back to \emph{*ṭhēḍḍ}, \textquotedblleft{}push\textquotedblright{}, a form CDIAL reconstructs (CDIAL 5512), which lists Hindi \emph{ṭhelā}, \textquotedblleft{}a push-cart\textquotedblright{}: the cart you push.
\end{cousinweb}

\begin{grammarlens}[title={What you've built}]
One more word for getting around.
\end{grammarlens}

\subsection*{Your turn}
\begin{itemize}
\raggedright
  \item \emph{Say it:} \emph{ṭhelā}
  \item \emph{Say it:} \emph{ṭhelā}, once more
  \item \emph{Say it:} \emph{ek ādmī ṭhele par sabzī bec rahā thā.} — A man was selling vegetables from a handcart.
  \item \emph{Recall:} say the Hindi for a roadside eatery, then the Hindi for a sweet-maker, then say \emph{ṭhelā} again
\end{itemize}

\begin{tcolorbox}[breakable,colback=teal!4,colframe=teal!35!black,title={Before you move on}]
What does \hi{ठेला} mean? (\textquotedblleft{}A handcart\textquotedblright{}.) Say it once more.
\end{tcolorbox}
\section[\texorpdfstring{\hi{फेरीवाला} (pherīvālā)}{फेरीवाला (pherīvālā)}]{\hi{फेरीवाला} (pherīvālā) — a hawker, a street seller}

\label{lesson:HI-C424-pherivala}



\begin{quote}
Before the new one: say the Hindi for a sweet-maker, then the Hindi for a handcart.
\end{quote}

\subsection*{You'll want to know: \hi{फेरीवाला}}
\textbf{\hi{फेरीवाला}} — \emph{pherīvālā} — \textquotedblleft{}a hawker, a street seller\textquotedblright{}.

\begin{quote}
\hi{फेरीवाला} \hi{हर} \hi{सुबह} \hi{गली} \hi{में} \hi{आता} \hi{है।}
\end{quote}

\begin{quote}
\emph{pherīvālā har subah galī meṁ ātā hai.}
\end{quote}

\begin{quote}
The hawker comes to the lane every morning.
\end{quote}

Masculine.

\begin{cousinweb}
\textbf{\hi{फेरीवाला}} (\emph{pherīvālā}) is \emph{pherī}, \textquotedblleft{}a round\textquotedblright{}, with -\hi{वाला}: the man who makes the rounds of the lanes. \emph{pherī} belongs with \emph{*phirati}, \textquotedblleft{}turns, wanders\textquotedblright{}, a form CDIAL reconstructs (CDIAL 9078), which lists Hindi \emph{pher}, \textquotedblleft{}a turn\textquotedblright{}.
\end{cousinweb}

\begin{grammarlens}[title={What you've built}]
One more word for getting around.
\end{grammarlens}

\subsection*{Your turn}
\begin{itemize}
\raggedright
  \item \emph{Say it:} \emph{pherīvālā}
  \item \emph{Say it:} \emph{pherīvālā}, once more
  \item \emph{Say it:} \emph{pherīvālā har subah galī meṁ ātā hai.} — The hawker comes to the lane every morning.
  \item \emph{Recall:} say the Hindi for a sweet-maker, then the Hindi for a handcart, then say \emph{pherīvālā} again
\end{itemize}

\begin{tcolorbox}[breakable,colback=teal!4,colframe=teal!35!black,title={Before you move on}]
What does \hi{फेरीवाला} mean? (\textquotedblleft{}A hawker, a street seller\textquotedblright{}.) Say it once more.
\end{tcolorbox}
\section[\texorpdfstring{\hi{पान} (pān)}{पान (pān)}]{\hi{पान} (pān) — a betel leaf, paan}

\label{lesson:HI-C424-pan}



\begin{quote}
Before the new one: say the Hindi for a handcart, then the Hindi for a hawker.
\end{quote}

\subsection*{You'll want to know: \hi{पान}}
\textbf{\hi{पान}} — \emph{pān} — \textquotedblleft{}a betel leaf, paan\textquotedblright{}.

\begin{quote}
\hi{खाने} \hi{के} \hi{बाद} \hi{दादा} \hi{जी} \hi{पान} \hi{खाते} \hi{हैं।}
\end{quote}

\begin{quote}
\emph{khāne ke bād dādā jī pān khāte haiṁ.}
\end{quote}

\begin{quote}
Grandfather has a paan after his meal.
\end{quote}

Masculine.

\begin{cousinweb}
\textbf{\hi{पान}} (\emph{pān}) is a tadbhava of Sanskrit \emph{parṇa}, \textquotedblleft{}a leaf\textquotedblright{} (CDIAL 7918), which lists Hindi \emph{pān}, \textquotedblleft{}a betel leaf\textquotedblright{}: the leaf above all leaves.
\end{cousinweb}

\begin{grammarlens}[title={What you've built}]
One more word for getting around.
\end{grammarlens}

\subsection*{Your turn}
\begin{itemize}
\raggedright
  \item \emph{Say it:} \emph{pān}
  \item \emph{Say it:} \emph{pān}, once more
  \item \emph{Say it:} \emph{khāne ke bād dādā jī pān khāte haiṁ.} — Grandfather has a paan after his meal.
  \item \emph{Recall:} say the Hindi for a handcart, then the Hindi for a hawker, then say \emph{pān} again
\end{itemize}

\begin{tcolorbox}[breakable,colback=teal!4,colframe=teal!35!black,title={Before you move on}]
What does \hi{पान} mean? (\textquotedblleft{}A betel leaf, paan\textquotedblright{}.) Say it once more.
\end{tcolorbox}
